\BeleskeID{08-s047}
% element: 08-s047:e04
% element: 08-s047:notes
Polazi se od početnog napona kondenzatora i opšteg rešenja. U trenutku početka prelazne pojave eksponencijalni faktor postaje jedan, pa napon određuje zbir konstanti. Drugi odnos dobija se poređenjem početnog nagiba iz diferencijalne jednačine i iz izvoda opšteg rešenja.

Član $K_2$ je konstantan i zato mu je izvod nula. Diferenciranje eksponencijalnog člana daje faktor $-1/\tau$. Ova dva odnosa određuju konstante: $K_2$ je vezana za konstantnu pobudu, a $K_1$ za početno stanje. Ne uvodi se novo nezavisno fizičko stanje.
